% =====================================================================
%  Chapter 24 — \idxw{Dust}, \idxw{Night} and the Engineering Envelope
% =====================================================================

\chapter{Dust, Night and the Engineering Envelope}
\label{ch:envelope}

The lunar concept has two objections that have been attached to it since
the beginning: the dust\idx{Dust} and the night. The dust\idx{Dust} objection has moved
furthest, because a mitigation technology that was a proposal when the
concept was written has now been demonstrated on the Moon. The night\idx{Night}
objection has not moved at all, because it is arithmetic.

\section{Dust, and the demonstration that changed its status}

The \idxw{Electrodynamic Dust Shield} lifts and clears regolith from surfaces using
a patterned electrode array generating a non-uniform electric field. It was
demonstrated on the lunar surface during Firefly Aerospace's Blue Ghost
Mission 1 in 2025, raising its technology readiness level from 5 to 7, with
a camera-lens clearing efficiency estimated at 96.6 per cent. NASA's current
dust\idx{Dust} mitigation catalogue lists it at that level for the Moon.

Two things should be said about it, and the second is the important one.

The first is that the dust objection has changed character. It is no longer
an open question whether the dust can be moved off an optic. It is a
question about how long an optic can be kept clean, on a surface where the
optics cannot be serviced for years, with a mitigation that NASA's own
roadmap cautions ``should not be the only tool in our toolbox''.

The second is that a shield is a mitigation for dust \emph{on a surface}.
Nothing in the published work shows that it can protect the primary mirror of
an interferometric aperture, or the baffle edges of a kilometre-scale array
spread over three kilometres of ground, or the coatings on a delay line that
has to hold a path difference to a fraction of a wavelength. The
demonstration was a camera lens on a lander.

\begin{seamnote}{The distinction the proposal has to live with}
A camera lens on a lander is a flat, exposed, near-surface, replaceable
optical element. A kilometre-scale array is a set of apertures on anchor
points, each with its own baffle, and the array is not replaceable. The
dust mitigation that has been demonstrated addresses the first problem and
not the second. The book does not claim otherwise; it claims that the
engineering envelope is closed enough to be worth designing against, and
that the residual dust risk is one of the things that would have to be
managed by the baseline selection rule rather than by cleaning.
\end{seamnote}

\section{The night, which is arithmetic}

The lunar night\idx{Night} at a polar site runs for about 328 hours, or thirteen and a
two-thirds days. Over that period a thermally exposed instrument drifts
towards the temperature of its surroundings, and the published engineering
figures for a thermally controlled payload at the south pole give survival
temperatures of about 100 kelvin and a seven kilowatt-hour battery to carry
a spectrometer through the night.

Three consequences follow for the proposal of Chapter 25.

The optical channel's baffle has to be held below about 250 kelvin for the
background arithmetic of Chapter 23 to work, and that is an insulation
problem rather than a cryogenics problem. The mid-infrared channel, which
needs 40 kelvin, would need the permanently shadowed region and therefore
the full power architecture. This is precisely why the proposed architecture
has no mid-infrared channel: the night\idx{Night} is what makes the cryogenic option
expensive, and the optical option is cheap by comparison.

The instrument cannot be pointed at the Sun during the night in any case,
because the Sun is below the horizon at a polar site for the whole of it.
The operational consequence is that the night is not a period of reduced
service but a period of no service, and that the campaign scheduling has to
account for it.

And the station's power has to come from somewhere that is not the
illuminated surface, or from storage that lasts the night. The published
answer --- regenerative fuel cells combined with vertical solar arrays,
delivering 10 kilowatts to permanently shadowed regions --- is a
kilowatt-class architecture for an instrument that does not need kilowatts.

\begin{table}[!htbp]
\centering
\small
\caption{The engineering envelope, item by item, with the status each item
has reached.}
\label{tab:envelope}
\begin{tabularx}{\textwidth}{@{}P{1.37in}P{1.19in}Y@{}}
\toprule
\rowcolor{panel}
\textbf{Item} & \textbf{Status} & \textbf{What the book claims} \\
\midrule
Atmospheric seeing & solved & no atmosphere at all \\
\rowcolor{panel}
Baseline stability & solved & a rigid, airless, slowly rotating body \\
Dust & mitigated & a flight demonstration at TRL 7; the long-term case for
a kilometre-scale array is open \\
\rowcolor{panel}
Night survival & arithmetic & 328 hours, 100 K survival, a 7 kWh battery \\
Power & published & 10 kW to permanently shadowed regions \\
\rowcolor{panel}
Thermal background & closed for optical & the baffle at 250 K; the
mid-infrared channel loses to L1 \\
Observing window & quantified & 6.83 days for the orthogonal orientation \\
\bottomrule
\end{tabularx}
\end{table}

\section{What the envelope permits}

The envelope in \tabref{tab:envelope} permits a specific instrument and
forbids others. It permits a visible-light interferometer with a thermally
controlled baffle, a surveyed array of anchor points, a service that operates
in daylight near the pole, and a campaign structure that uses the slow
rotation of the sky.

It forbids a mid-infrared channel, a cryogenic payload in an illuminated
site, and a service that depends on a continuously available sky. And it
imposes the constraint that the station's duty cycle is set by the
twenty-eight-day lunar cycle rather than by the observing target.

That last point is worth stating as a number, because it is the number that
the campaign rate in the next chapter depends on. The station can observe
for roughly half of each synodic month, at the pole, and it can observe
objects that are near the Sun at all, and it cannot observe through the
night. The net effect on the annual service rate is a factor that the
book's generator folds into its overhead allowance.

\begin{keytakeaways}
\item Dust mitigation has been demonstrated on the Moon at TRL 7 with a
96.6 per cent camera-lens clearing efficiency, but it has not been
demonstrated for a kilometre-scale array.
\item The night runs for 328 hours, with survival temperatures near 100
kelvin and a 7 kilowatt-hour battery.
\item The envelope permits a visible-light instrument with a baffle at about
250 kelvin and forbids a mid-infrared channel in an illuminated site.
\item The duty cycle is set by the lunar cycle rather than by the target,
and it is the single largest term in the overhead allowance.
\end{keytakeaways}
